
\documentclass{article}
\usepackage{colortbl}
\usepackage{makecell}
\usepackage{multirow}
\usepackage{supertabular}

\begin{document}

\newcounter{utterance}

\centering \large Interaction Transcript for game `mmlu\_pro', experiment `mmlu\_pro\_biology', episode 481 with reflexion{-}Qwen3.5{-}4B.
\vspace{24pt}

{ \footnotesize  \setcounter{utterance}{1}
\setlength{\tabcolsep}{0pt}
\begin{supertabular}{c@{$\;$}|p{.15\linewidth}@{}p{.15\linewidth}p{.15\linewidth}p{.15\linewidth}p{.15\linewidth}p{.15\linewidth}}
    \# & \multicolumn{2}{c}{Player} && \multicolumn{2}{c}{Game Master} \\
    \hline

    \theutterance \stepcounter{utterance}  
    & & & \multicolumn{4}{p{0.6\linewidth}}{
        \cellcolor[rgb]{0.9,0.9,0.9}{
            \makecell[{{p{\linewidth}}}]{
                \texttt{\tiny{[P1$\langle$GM]}}
                \texttt{You have attempted this same task before and did not succeed. These are the plans you wrote after those attempts.} \\
\texttt{Trial \#1: New plan: I failed because I answered the final question immediately without verifying the biological calculations for the specific stages of meiosis requested. The question asks for chromosome and chromatid counts at four distinct points: spermatogonium (diploid, pre{-}replication), first polar body (haploid, post{-}meiosis I), second polar body (haploid, post{-}meiosis II), and secondary oocyte (haploid, pre{-}meiosis II). I likely guessed based on general knowledge of haploid/diploid states rather than tracking the chromatid count (which doubles after S{-}phase and halves after meiosis I/II). Next time, I will explicitly calculate the chromosome and chromatid numbers for each of the four specific stages before selecting the option, ensuring the chromatid counts reflect the cell cycle stage (e.g., spermatogonium has 20 chromosomes and 20 chromatids if in G2, or 20/40 depending on timing, but the options clarify the intended stage). I will carefully match the option to the precise definition of each cell type in the sequence.} \\
\\ 
\texttt{Trial \#2: New plan: I failed because I selected option (H) without verifying the specific chromatid counts for each stage, likely confusing the number of chromatids in the secondary oocyte (which should be 20 chromatids, as it is a haploid cell with replicated chromosomes) against the second polar body. The secondary oocyte is haploid (10 chromosomes) but has replicated chromosomes, meaning it contains 20 chromatids. The second polar body is also haploid with replicated chromosomes (10 chromosomes, 20 chromatids). Option (H) incorrectly lists 10 chromatids for the secondary oocyte. Next time, I will explicitly calculate the chromosome and chromatid numbers for all four stages: spermatogonium (20 chromosomes, 20 chromatids), first polar body (10 chromosomes, 20 chromatids), second polar body (10 chromosomes, 20 chromatids), and secondary oocyte (10 chromosomes, 20 chromatids), then select the option that matches these exact figures.} \\
\\ 
\\ 
\texttt{The following are multiple choice questions (with answers) about biology. Answer only with "The answer is (X)" where X is the correct letter choice.} \\
\\ 
\texttt{Question:} \\
\texttt{Which of the following represents an accurate statement concerning arthropods?} \\
\texttt{Options:} \\
\texttt{A. They possess an exoskeleton composed primarily of peptidoglycan.} \\
\texttt{B. They possess an open circulatory system with a dorsal heart.} \\
\texttt{C. They are members of a biologically unsuccessful phylum incapable of exploiting diverse habitats and nutrition sources.} \\
\texttt{D. They lack paired, jointed appendages.} \\
\texttt{E. N/A} \\
\texttt{F. N/A} \\
\texttt{G. N/A} \\
\texttt{H. N/A} \\
\texttt{I. N/A} \\
\texttt{J. N/A} \\
\texttt{The answer is (B)} \\
\\ 
\texttt{Question:} \\
\texttt{In a given population, 1 out of every 400 people has a cancer caused by a completely recessive allele, b. Assuming the population is in Hardy{-}Weinberg equilibrium, which of the following is the expected proportion of individuals who carry the b allele but are not expected to develop the cancer?} \\
\texttt{Options:} \\
\texttt{A. 19/400} \\
\texttt{B. 1/400} \\
\texttt{C. 40/400} \\
\texttt{D. 38/400} \\
\texttt{E. 2/400} \\
\texttt{F. 1/200} \\
\texttt{G. 20/400} \\
\texttt{H. 50/400} \\
\texttt{I. N/A} \\
\texttt{J. N/A} \\
\texttt{The answer is (D)} \\
\\ 
\texttt{Question:} \\
\texttt{A mutation in a bacterial enzyme changed a previously polar amino acid into a nonpolar amino acid. This amino acid was located at a site distant from the enzyme’s active site. How might this mutation alter the enzyme’s substrate specificity?} \\
\texttt{Options:} \\
\texttt{A. By changing the enzyme’s pH optimum} \\
\texttt{B. By changing the enzyme's molecular weight} \\
\texttt{C. An amino acid change away from the active site increases the enzyme's substrate specificity.} \\
\texttt{D. By changing the shape of the protein} \\
\texttt{E. By changing the enzyme's temperature optimum} \\
\texttt{F. By altering the enzyme's ability to be denatured} \\
\texttt{G. By changing the enzyme’s location in the cell} \\
\texttt{H. By changing the enzyme's color} \\
\texttt{I. An amino acid change away from the active site cannot alter the enzyme’s substrate specificity.} \\
\texttt{J. By altering the enzyme's rate of reaction} \\
\texttt{The answer is (D)} \\
\\ 
\texttt{Question:} \\
\texttt{Which of the following is not a way to form recombinant DNA?} \\
\texttt{Options:} \\
\texttt{A. Translation} \\
\texttt{B. Conjugation} \\
\texttt{C. Specialized transduction} \\
\texttt{D. Transformation} \\
\texttt{E. N/A} \\
\texttt{F. N/A} \\
\texttt{G. N/A} \\
\texttt{H. N/A} \\
\texttt{I. N/A} \\
\texttt{J. N/A} \\
\texttt{The answer is (A)} \\
\\ 
\texttt{Question:} \\
\texttt{Which of the following is not known to be involved in the control of cell division?} \\
\texttt{Options:} \\
\texttt{A. Microtubules} \\
\texttt{B. Checkpoints} \\
\texttt{C. DNA polymerase} \\
\texttt{D. Centrosomes} \\
\texttt{E. Cyclins} \\
\texttt{F. Mitochondria} \\
\texttt{G. Protein kinases} \\
\texttt{H. Fibroblast cells} \\
\texttt{I. N/A} \\
\texttt{J. N/A} \\
\texttt{The answer is (H)} \\
\\ 
\texttt{Question:} \\
\texttt{In an animal with a haploid number of 10, how many chromosomesare present in (a)aspermatogonium? (b)inthe first polar body? (c)inthe second polar body? (d)inthe secondaryoocyte? Assume that the animal is diploid.} \\
\texttt{Options:} \\
\texttt{A. 20 chromosomes, 30 chromatids, 10 chromosomes, 15 chromatids} \\
\texttt{B. 20 chromosomes, 40 chromatids, 20 chromosomes, 40 chromatids} \\
\texttt{C. 10 chromosomes, 10 chromatids, 20 chromosomes, 10 chromatids} \\
\texttt{D. 20 chromosomes, 10 chromatids, 20 chromosomes, 10 chromatids} \\
\texttt{E. 10 chromosomes, 5 chromatids, 10 chromosomes, 5 chromatids} \\
\texttt{F. 15 chromosomes, 15 chromatids, 15 chromosomes, 15 chromatids} \\
\texttt{G. 20 chromosomes, 20 chromatids, 10 chromosomes, 20 chromatids} \\
\texttt{H. 10 chromosomes, 20 chromatids, 10 chromosomes, 20 chromatids} \\
\texttt{I. 20 chromosomes, 10 chromatids, 10 chromosomes, 5 chromatids} \\
\texttt{J. 5 chromosomes, 10 chromatids, 5 chromosomes, 10 chromatids} \\
            }
        }
    }
     \\ \\

    \theutterance \stepcounter{utterance}  
    & \multicolumn{4}{p{0.6\linewidth}}{
        \cellcolor[rgb]{0.561,0.737,0.561}{
            \makecell[{{p{\linewidth}}}]{
                \texttt{\tiny{[P1$\rangle$GM]}}
                \texttt{The answer is (G)} \\
            }
        }
    }
    & & \\ \\

    \theutterance \stepcounter{utterance}  
    & & & \multicolumn{2}{p{0.3\linewidth}}{
        \cellcolor[rgb]{0.9,0.9,0.9}{
            \makecell[{{p{\linewidth}}}]{
                \texttt{\tiny{[GM$|$GM]}}
                \texttt{G} \\
            }
        }
    }
    & & \\ \\

    \theutterance \stepcounter{utterance}  
    & & & \multicolumn{2}{p{0.3\linewidth}}{
        \cellcolor[rgb]{0.9,0.9,0.9}{
            \makecell[{{p{\linewidth}}}]{
                \texttt{\tiny{[GM$|$GM]}}
                \texttt{G} \\
            }
        }
    }
    & & \\ \\

    \theutterance \stepcounter{utterance}  
    & & & \multicolumn{2}{p{0.3\linewidth}}{
        \cellcolor[rgb]{0.9,0.9,0.9}{
            \makecell[{{p{\linewidth}}}]{
                \texttt{\tiny{[GM$|$GM]}}
                \texttt{game\_result = WIN} \\
            }
        }
    }
    & & \\ \\

\end{supertabular}
}

\end{document}
